\begin{frame}
  \frametitle{Problem Descriptions}
  The results slides will solve the same benchmarks with the same discretizations as the original discrete-ordinates methods to demonstrate the removal of the ray effect via spherical harmonics. \\

  As a reminder, the ray effect is the persistence of flux in the solvers sweeping directions, and is particularly apparent in geometries with low scattering.
\end{frame}

\begin{frame}
  \frametitle{\glspl{arp1} Absorbing Medium Results}
  \begin{figure}[b]
    \centering
    \begin{subfigure}{0.45\textwidth}
      \centering
      \includegraphics[width=0.65\textwidth]{figures/absorbing/cell_fluxes_256_8.png}
      \subcaption{$N=256$}
    \end{subfigure}%
    \begin{subfigure}{0.45\textwidth}
      \centering
      \includegraphics[width=0.65\textwidth]{figures/absorbing/cell_fluxes_1024_8.png}
      \subcaption{$N=1024$}
    \end{subfigure}
    \begin{subfigure}{0.45\textwidth}
      \centering
      \includegraphics[width=0.65\textwidth]{figures/absorbing/cell_fluxes_4096_8.png}
      \subcaption{$N=4096$}
    \end{subfigure}%
    \begin{subfigure}{0.45\textwidth}
      \centering
      \includegraphics[width=0.65\textwidth]{figures/absorbing/cell_fluxes_16384_8.png}
      \subcaption{$N=16384$}
    \end{subfigure}
    \caption{Solution of absorbing \glspl{arp1} for varying $N$ with $S_{8}$.}
  \end{figure}
\end{frame}
